\scanpage{12}
% Pre-ruled music paper.
$\Lambda^k(\Lambda^i\Phi_{*k})$ et $\Lambda^k(\Lambda^j\Psi_{*k})$ sont spéciaux (bien sûr, ils seront $\overline{\mathrm m}$ \unsure{strts|très} spéciaux).

Leurs composantes sont évidemment libres. De plus, comme $\Lambda^i\Phi_{*k}$ est homotope à $0$, de $\overline{\mathrm m}$ $\Lambda^k(\Lambda^i\Phi_*)$, dont les $\pi_i$ sont donc nuls et à fortiori libres de t.f. :
\margin{?}
\hfill $]$

\uline{Remarque}. Soit $L_{\bullet}$ $k$-module $\frac12$ simplicial \deleted{\unsure{str}} spécial, alors $L_{\bullet}$ est connu à isom près, i.e les $\mu_i=\operatorname{rg}M_i$ $(i\geq1)$ et $\nu_j=\operatorname{rg}N_j$ \inserted{$j\geq0$,} sont connus, quand on connait la \deleted{\unsure{valeurs}} \inserted{fonction} $n\longmapsto\operatorname{rg}_k L_n=f_{L_{\bullet}}(n)$, \deleted{En effet} \uline{et} les $\operatorname{rg}_k\deleted{\unsure{H}}\pi_j(L_{\bullet})$ $(=\operatorname{rg}_k N_j=\nu_j)$. En effet, on aura
\[
\sum_{i\geq1}\mu_i\underset{\substack{\Vert\\P_i(n+1)}}{\binom{n+1}i}
=f_{L_{\bullet}}(n)-\sum_{j\geq0}\nu_j\binom n j
\qquad(n\ \text{grand})
\]
\deleted{Or on sait} i.e.
\[
f_{L_{\bullet}}(n-1)-\sum_{j\geq0}P_j(n-1\deleted{\unsure{j}})=\sum_{i\geq1}\mu_i P_i(n)
\]
or on sait que les \inserted{pol $P_i(n)$} \deleted{\unsure{$\binom n i$}} sont \deleted{des} lin. \unsure{indépendts|indépendants} --- et $\overline{\mathrm m}$ forment une base du $\mathbb Z$-module des polynômes à \inserted{coeff. entiers et} valeurs entières sur les entiers.

Donc la connaissance des $\sum_{i\geq1}\mu_iP_i$ pour $n$ grand implique
% archive sheet number: 11
